% Figure: a gravity retaining wall and its three stability checks. The wall
% (a trapezoidal masonry block) retains soil on its left. The active earth
% pressure rises triangularly with depth, resultant Pa at one-third height.
% The wall weight W acts down through the block centroid. Stability is read
% from three checks: overturning about the toe, sliding along the base, and
% bearing under the base.
\begin{figure}[htbp]
\centering
\begin{tikzpicture}[>=latex, x=1.05cm, y=1.05cm]
  % wall cross-section: trapezoid, wider at the base (gravity wall)
  % heel-side (retained soil) on the left, toe on the right
  \coordinate (BL) at (0,0);     % base left (heel)
  \coordinate (BR) at (3.0,0);   % base right (toe)
  \coordinate (TL) at (1.0,4.0); % top left
  \coordinate (TR) at (2.0,4.0); % top right
  \draw[engblue, very thick, fill=engblue!10] (BL) -- (BR) -- (TR) -- (TL) -- cycle;
  % retained soil to the left (hatched ground)
  \fill[enggray!18] (BL) -- (TL) -- (-2.6,4.0) -- (-2.6,0) -- cycle;
  \foreach \y in {0.4,0.9,...,3.9} {
    \draw[enggray, thin] (-2.6,\y) -- (-2.45,\y+0.18);
  }
  % ground surface line
  \draw[enggray, thick] (-2.6,4.0) -- (TL);
  % base ground line
  \draw[engblack, thick] (-0.4,0) -- (3.6,0);
  % active earth pressure triangle on the back face (heel side).
  % use a representative vertical face at x=0 for the pressure diagram;
  % horizontal arrows grow with depth, pointing into the wall (rightward)
  \def\xheel{0}
  \foreach \k in {1,2,3,4,5} {
    \pgfmathsetmacro{\yy}{\k*0.7}
    \pgfmathsetmacro{\len}{0.95*(4.0-\yy)/4.0}
    \draw[->, engred, thin] ({\xheel-\len},\yy) -- (\xheel,\yy);
  }
  % resultant active thrust Pa at one-third height
  \draw[->, very thick, engred] (-1.5,{4.0/3}) -- (\xheel,{4.0/3});
  \node[engred, font=\scriptsize, anchor=east] at (-1.55,{4.0/3}) {$P_a$};
  \node[engred, font=\scriptsize, anchor=west] at (0.05,{4.0/3 - 0.35}) {at $h/3$};
  % wall weight at centroid (approx centroid of trapezoid)
  \coordinate (G) at (1.45,1.7);
  \fill[engblack] (G) circle (1.4pt);
  \draw[->, very thick, engblack] (G) -- ++(0,-1.5);
  \node[font=\scriptsize, anchor=west] at ($(G)+(0.08,-0.9)$) {$W$};
  % base reaction resultant (vertical) near the toe, and friction along base
  \draw[->, thick, englime!70!black] (2.2,0) -- ++(0,1.2);
  \node[font=\scriptsize, anchor=west] at (2.25,0.8) {$N$};
  \draw[->, thick, englime!70!black] (2.2,0) -- ++(1.0,0);
  \node[font=\scriptsize, anchor=south] at (2.9,0.05) {$F=\mu N$};
  % toe marker (overturning is taken about this point)
  \fill[engamber] (BR) circle (2.0pt);
  \node[engamber, font=\scriptsize, anchor=north west] at ($(BR)+(0.04,-0.05)$) {toe $O$};
  % height dimension
  \draw[<->, enggray, thin] (-2.95,0) -- (-2.95,4.0);
  \node[enggray, font=\scriptsize, anchor=east] at (-2.95,2.0) {$h$};
  % base width dimension
  \draw[<->, enggray, thin] (0,-0.55) -- (3.0,-0.55);
  \node[enggray, font=\scriptsize, anchor=north] at (1.5,-0.55) {base width $B$};
\end{tikzpicture}
\caption[Gravity retaining wall and its three stability checks]{A gravity
retaining wall, held in place by its own weight against the earth pressure
of the soil it retains. The active pressure grows linearly with depth and
sums to a horizontal thrust $P_a=\tfrac12 K_a \gamma h^2$ per unit length
of wall, acting at one-third the height above the base. The wall weight $W$
acts down through the block centroid. Three checks decide whether the wall
stands. Overturning compares the restoring moment of $W$ about the toe $O$
against the overturning moment of $P_a$ about the same point. Sliding
compares the base friction $\mu N$ against the thrust $P_a$. Bearing
compares the peak pressure under the base against the soil's allowable
bearing capacity. A gravity wall earns its stability from weight and base
width, the two quantities the three checks trade against the retained
height.}
\label{fig:vol03ch02:retaining-wall}
\end{figure}
